\label{sec:introduccion}

El transporte de electrones en gases está determinado por la interacción entre las partículas cargadas, el campo eléctrico aplicado y las especies neutras del medio. Las colisiones elásticas e inelásticas modifican la energía y la dirección de los electrones y, en conjunto, dan origen a magnitudes macroscópicas como la movilidad, la velocidad de deriva, los coeficientes de difusión y las tasas de ionización. El conocimiento de estos parámetros es necesario para caracterizar descargas gaseosas y evaluar modelos cinéticos de plasmas fuera del equilibrio \cite{GonzalezMagana2024,Jovanovic2009}.\\

En los gases nobles, la probabilidad de dispersión de electrones de baja energía no presenta un comportamiento monótono. El efecto Ramsauer--Townsend consiste en la aparición de un mínimo en la sección eficaz de dispersión para un intervalo particular de energía electrónica, de modo que el gas se vuelve relativamente transparente al paso de los electrones \cite{Ramsauer1921}. Este fenómeno puede observarse en un tiratrón de xenón 2D21 mediante la medición de las corrientes de rejilla y ánodo mientras se modifica la energía de los electrones emitidos por el cátodo \cite{Kukolich1968}. En la primera parte del trabajo se utilizó este dispositivo para estudiar la respuesta eléctrica asociada al mínimo Ramsauer--Townsend y comparar la adquisición manual con un sistema automatizado.\\

Una descripción complementaria del transporte se obtiene mediante los experimentos de enjambres, en los que un conjunto de electrones atraviesa un gas bajo la acción de un campo aproximadamente uniforme. Sus coeficientes de transporte dependen principalmente del campo eléctrico reducido $E/N$, donde $E$ es la magnitud del campo y $N$ la densidad numérica del gas \cite{GonzalezMagana2019}. En la técnica de Townsend pulsada, un pulso ultravioleta libera electrones desde el cátodo y el movimiento posterior de electrones e iones produce una corriente transitoria. Debido a la diferencia entre sus movilidades, la componente electrónica puede separarse temporalmente y utilizarse para determinar la velocidad de deriva \cite{Jovanovic2009}.

Las mezclas de argón y nitrógeno son particularmente sensibles a la competencia entre la dispersión elástica del Ar, relacionada con el mínimo Ramsauer--Townsend, y los canales de excitación del $\mathrm{N_2}$. Para concentraciones pequeñas de nitrógeno, esta interacción puede producir conductividad diferencial negativa, es decir, una disminución de la velocidad de deriva al aumentar $E/N$ dentro de cierto intervalo \cite{Jovanovic2009}. Por ello, en la segunda parte se midió la corriente electrónica transitoria en una mezcla de \SI{1}{\percent} de $\mathrm{N_2}$ y \SI{99}{\percent} de Ar, y se obtuvo la velocidad de deriva en función de $E/N$. En conjunto, ambas experiencias permitieron examinar cómo las colisiones electrón--átomo se manifiestan tanto en las corrientes de un tubo gaseoso como en el transporte colectivo de un enjambre electrónico.